\documentclass{article}
\usepackage{amsmath,amssymb}
\usepackage{cite}

% Needs LaTeX 2019 or newer. Writes the .bib out at compile time
% \begin{filecontents}[overwrite]{references.bib}
% @article{author2026,
%   title = {Title of paper},
%   author = {Author, Firstname and Coauthor, Secondname},
%   journal = {Journal of Mathematical Data Science},
%   volume = {10},
%   number = {2},
%   pages = {123--145},
%   year = {2026},
%   doi = {10.1000/182}
% }
% \end{filecontents}

\title{Arctan-skewing: a new skewing scheme on the hypertorus}

% Speaker only, with affiliation and optional ORCID
\author{
  Sophia Loizidou\\
  \textit{Department of Mathematics, University of Luxembourg, Luxembourg}\\
  ORCID: 0009-0006-1372-7032
}
\date{}

\begin{document}

\maketitle

\begin{abstract}
In this talk we propose and study arctan-skewing, a general approach to skew any symmetric distribution on the unit hypertorus. 
Arctan-skewing retains the main advantages of the popular sine-skewing mechanism, but improves on it by overcoming its two major limitations. 
These are a condition on the skewness parameter and a singular Fisher Information matrix in the vicinity of symmetry when combined with well-known symmetric models from the literature such as the circular von Mises or bivariate Cosine models. 
Arctan-skewing thus represents a compelling alternative to model skewed data on the circle and torus, notably from the field of protein bioinformatics. 
We present its stochastic properties and a maximum likelihood estimation procedure, and devise optimal tests for symmetry against our new class of skewed alternatives. 
Finally we will show that arctan-skewed models better fit protein data compared to their sine-skewed counterparts.

Joint work with Christophe Ley.
\end{abstract}

\bibliographystyle{plain}
\bibliography{references}

\end{document}